Write $S = \bK[x_1,\ldots,x_n]$ and $I_n = (e_1,\ldots,e_n)$ for the ideal of $S$ generated by the elementary symmetric polynomials. 
A classical result of Borel \cite{Borel1953} states that for any field $\bK,$ the cohomology ring of the complete flag variety $\fl_n$ has a presentation
\[
H^*(\fl_n;\bK)\cong R_n,
\]
where $R_n = S/I_n$ is the \emph{coinvariant algebra}. 
For any Schubert variety $X_w\subseteq \fl_n$, the restriction map $H^*(\fl_n;\bK)\twoheadrightarrow H^*(X_w;\bK)$ is surjective and its kernel is the ideal $J_w\subseteq R_n$ linearly spanned by the Schubert polynomials $\fraks_v$ with $v\not\leq w.$ 
Hence, we have a presentation of the cohomology ring of a Schubert variety as $H^*(X_w;\bK)\cong R_n/J_w.$

While the Schubert polynomials $\fraks_v$ for $v\not\leq w$ give a convenient $\bK$-linear spanning set of $J_w$ (modulo $I_n$), it is still an open problem to give a minimal generating set of the ideal $J_w$. Toward this goal, Reiner, Woo, and Yong \cite{ReinerWooYong} found a relatively small generating set of $J_w$ consisting of Schubert polynomials indexed by bigrassmannian permutations.
By work of St. Dizier-Yong \cite{stdizier-yong}, $J_w$ can be written as a sum of \emph{basic ideals}, and each basic ideal is minimally generated by a subset of Reiner--Woo--Yong generators.

On the other hand, Gasharov and Reiner gave a particularly small generating set of order $n^2$ for $J_w$ in the case that $w$ is \emph{defined by inclusions}, meaning that its associated Schubert variety $X_w$ is cut out of $\fl_n$ by inclusion conditions on the parts of its flag. 
Gasharov and Reiner characterize $w\in\sn$ defined by inclusions as those avoiding the patterns $4231,35142,42513,$ and $351624$ and show that $J_w$ is generated by the elementary symmetric polynomials in initial and final blocks of the variables. 
Ding \cite{ding01} has also given a family of Schubert varieties in partial flag varieties, each of whose cohomology ideal has a generating set with up to $n$ generators. 
Note, while there is some overlap, Ding's class does not contain the class of Schubert varieties discussed in this paper and vice versa. 

We fix notation before stating our results. 
Given a set of permutations $A$, denote by $S_n(A)$ the set of permutations in $S_n$ that avoid all patterns in $A$. 
Let
\begin{align*}
\calP &\coloneqq \{ 2143,\ 2341,\ 2413,\ 2431,\ 4123,\ 4213,\ 4231 \}.
\end{align*}
For $w\in S_n$, we sometimes say \emph{$w$ is $\calP$-avoiding} to mean $w\in S_n(\calP)$. 
For a simplicial complex $\Delta$ on a vertex set $[n]:=\{1,\ldots,n\},$ let $I_\Delta\subseteq S$ be it's \emph{Stanley-Reisner ideal}, which is generated by the squarefree monomials $x_G$ determined by the nonfaces $G\subseteq [n]$ of $\Delta.$ The corresponding quotient ring $R_{\Delta} = S/I_\Delta$ is called the \emph{Stanley-Reisner ring} of $\Delta$. 


In this article, we introduce a family of simplicial complexes called \emph{doubly pruned} characterized by their minimal nonfaces being initial or final subsets of their support, see Definition~\ref{def: doubly pruned simplicial complexes} for the precise description. To each $w\in S_n(\calP)$, we attach a doubly pruned simplicial complex $\Delta_w$ and prove that the ideal $J_w\subseteq R_n$ is generated by square-free monomials. In fact, we prove these are exactly the Schubert varieties whose cohomology ring has this property. In our first main theorem, we prove that $J_w$ being generated by squarefree monomials in $R_n$ is equivalent to the pattern-avoidance condition above and to a set of inclusion conditions on the Schubert variety.

\begin{restatable}{thmintro}{MainA}
\label{thm:main-A}
    For $w\in S_n$, the following are equivalent:
    \begin{enumerate}%[(1)]
    \renewcommand{\theenumi}{\thethmintro\arabic{enumi}}
    \renewcommand{\labelenumi}{(\arabic{enumi})}
    \item\label[condition]{cond:A1} (Algebraic) The ideal $J_w$ is generated by squarefree monomials in $R_n$. 
    \item\label[condition]{cond:A2} (Combinatorial) $w\in\sn(\calP)$.
    \item\label[condition]{cond:A3}  (Geometric) The Schubert variety $X_w$ is defined by the following set \footnote{For $w=\id$, $X_\id$ is a single point defined by $V_k=\bC^k$ for all $k$.} of inclusions for some $i\leq j$
        \begin{itemize}
            \item $V_1 = \bC^1$,\dots, $V_{i-1} = \bC^{i-1}$,
            \item $\bC^i\subseteq V_{w^{-1}(i)}$,
            \item $V_{w^{-1}(j)-1}\subseteq \bC^{j-1}$,
            \item $V_{j} =\bC^{j}$,\dots, $V_{n-1} = \bC^{n-1}$
        \end{itemize}
    \end{enumerate}
\end{restatable}

Thus, if any of the equivalent conditions in \Cref{thm:main-A} holds, then $H^*(X_w;\bK)\cong S/J_w$ where $J_w$ is generated by squarefree monomials modulo $I_n$. Equivalently, $J_w = I_n +I_\Delta$ for some simplicial complex $\Delta$ on $[n]$. In this case, the set of generators of $J_w$ is particularly simple to describe and has size at most $n$.

In Section~\ref{sec: intro to doubly pruned simplicial complexes}, we give equivalent combinatorial conditions for a simplicial complex to be doubly pruned. Furthermore, the collection of doubly pruned complexes admits nice poset-theoretic properties. For example, in the assignment $w\in S_n(\calP)$ mapping to $\Delta_w$ gives an injection of the subposet of Bruhat order on $S_n(\calP)$ into the Boolean lattice, as we prove in \Cref{prop:Inclusion}. The image of this map of posets is properly contained in the collection of doubly pruned complexes. The subposet of doubly pruned complexes forms a sublattice of the Boolean lattice, as we prove in \Cref{cor: lattice}.





We extend \Cref{thm:main-A} to the setting of unions of Schubert varieties $X_\calW = \bigcup_{w\in \calW}X_w$. Since containment of Schubert varieties is governed by strong Bruhat order on their corresponding permutations, unions of Schubert varieties are indexed by antichain $\calW\subseteq S_n$ under Bruhat order. Set $J_\calW$ to be kernel of the surjective map $R_n = H^*(\fl_n;\bK)\to H^*(X_\calW;\bK)$. By standard arguments using dual bases, $J_\calW = \bigcap_{w\in \calW}J_w$. Further, set
\[
\Delta_\calW := \bigcup_{w\in\calW}\Delta_w.
\]
A basic property of Stanley--Reisner ideals is that the intersection of a collection of Stanley--Reisner ideals is equal to the Stanley--Reisner ideal of the union of the complexes, see~\cite{MillerSturmfels2005}. Thus, 
\[
I_{\Delta_\calW} = \bigcap_{w\in\calW}I_{\Delta_w}\subseteq S.
\]

In view of \Cref{thm:main-A}, one might expect that for $\calW\subseteq S_n(\calP)$,
\begin{equation}\label{eq: hopes and dreams}
J_\calW =\bigcap_{w\in \calW} J_w= \bigcap_{w\in \calW} (I_n + I_{\Delta_w}) =  I_n + \bigcap_{w\in \calW} I_{\Delta_w}= I_n + I_{\Delta_\calW},
\end{equation}
and thus $J_\calW$ is generated by squarefree monomials (modulo $I_n$). However, the third equality in \eqref{eq: hopes and dreams} does not always hold, even in the case of $|\calW|=2$, as the next example shows.

\begin{example} 
Let $n = 4$ and the commuting simple reflections $u = s_1 = 2134$ and $v = s_3 = 1243,$ so that $X_u$ and $X_v$ are projective curves meeting at $X_\id.$ Both lie in $\mathrm{S}_4(\calP),$ so $J_u$ and $J_v$ are generated by squarefree monomials, while $J_\calW$ is not. Indeed, since $s_2\not\leq s_1$ and $s_2\not\leq s_3,$ the class $\fraks_{s_2} = x_1 + x_2$ lies in $J_\calW,$ while no single variable is in $J_\calW$, as the reader can check. The failure of \eqref{eq: hopes and dreams} can then be seen as follows. The complex $\Delta_\calW$ is four isolated points and $S/(I_{\Delta_\calW}+I_4)$ has Hilbert series $1+3t,$ whereas for $X_\calW,$ we have two curves meeting at a point, which has cells indexed by $\id, s_1,$ and $s_3,$ and thus has Poincar\'e polynomial $1+2t.$  
\end{example}

The failure of \eqref{eq: hopes and dreams} is due to the fact that, in general, 
\[
I_n + (I_{\Delta_u}\cap I_{\Delta_v}) \subsetneq (I_n + I_{\Delta_u}) \cap (I_n + I_{\Delta_v}).
\]
We translate this failure into homological terms, as follows, and prove that it is equivalent to a topological condition on the associated complexes. Since $I_{\Delta_u}\cap I_{\Delta_v} = I _{\Delta_u\cup\Delta_v}$ and $I_{\Delta_u}+I_{\Delta_v} = I_{\Delta_u\cap\Delta_v},$ by lifting to the face rings, one may construct the Mayer-Vietoris sequence
\[
0\longrightarrow
S/I_{\Delta_u\cup\Delta_v}
\longrightarrow
S/I_{\Delta_u}\oplus S/I_{\Delta_v}
\longrightarrow
S/I_{\Delta_u\cap\Delta_v}
\longrightarrow
0.
\]
The cohomology rings of $X_u,X_v,$ and their union and intersection fit into an analogous short exact sequence, since Schubert varieties have no odd degree cohomology.
However, passing from face rings to cohomology rings means tensoring with $R_n = S/I_n,$ which is only right exact. The failure above is exactly the cokernel of the induced map on $\Tor_1^S(R_n,-)$, which we explain in \Cref{subsec:Prune_Equality}.

Our second main theorem (\Cref{thm:main-B}) shows that this failure is the only obstruction. More importantly, it shows that the only time that \Cref{eq: hopes and dreams} holds is for a Bruhat antichain $\calW\subseteq\sn$ such that $J_\calW$ is generated by squarefree monomials in $R_n.$ We also give a combinatorial-geometric interpretation.  
In order to prove this result (\Cref{sec: proofs of the main theorems}), we use Koszul resolutions to show that the map induced by $\Tor_1^S(R_n,-)$ is surjective if and only if the complexes $\Delta_w$ are glued together in a way reminiscent of a shelling on simplicial complexes (\Cref{def:prunable}).

\begin{restatable}[]{thmintro}{MainB}
\label{thm:main-B}
Let \(\mathcal W\subseteq S_n\) be a strong-Bruhat antichain. Then the following are equivalent.

\begin{enumerate}[(1)]
    \renewcommand{\theenumi}{\thethmintro\arabic{enumi}}
    \renewcommand{\labelenumi}{(\arabic{enumi})}
\item\label[condition]{cond:B1} (Algebraic) The ideal $J_{\mathcal W}$ is generated by squarefree monomials in \(R_n\) where
\[ J_{\mathcal W}:=\bigcap_{w\in\mathcal W} J_w\subseteq R_n. \]

\item\label[condition]{cond:B2} (Cohomological) We have \(\mathcal W\subseteq S_n(\calP)\) and 
\[ \bigcap_{w\in\mathcal W}(I_n+I_{\Delta_w}) = I_n+\bigcap_{w\in\mathcal W}I_{\Delta_w}\]
as ideals in $S=\bK[x_1,\dots,x_n]$.

\item\label[condition]{cond:B3} (Combinatorial) We have $\mathcal{W}\subseteq S_n(\mathcal{P})$, and the family of simplicial complexes $\{\Delta_w:w\in \mathcal{W}\}$ is \textbf{prunable}, which we define below in Definition~\ref{def:prunable}.\label[(3)]{item:3}
\end{enumerate}
\end{restatable}

%  Every element satisfies one of k Theorem A style inclusions
% There is an ordering of them such that 

The topological property we need in \eqref{item:3} above is as follows. Note that it is \textbf{not} a purely topological property of the union $\bigcup_{w\in \mathcal W}\Delta_w$ but rather depends on how the $\Delta_w$ are glued together.

\begin{definition}\label{def:prunable}
    Given a Bruhat antichain $\mathcal{W}\subseteq S_n(\mathcal{P})$, we say $\{\Delta_w:w\in \mathcal{W}\}$ is \textbf{prunable} if \(\mathcal W\) admits an ordering $w_1,\ldots,w_k$ such that the following two conditions hold for every \(2\le r\le k\), where
    \[\Omega_r\coloneqq\Delta_{w_r}\cap(\Delta_{w_1}\cup\cdots\cup\Delta_{w_{r-1}}).\]
    \begin{enumerate}
    \item[(a)]\(\Omega_r\) is pure and $\dim\Omega_r = \dim\Delta_{w_r}$.
    \item[(b)] For every nonempty subset \(T\subseteq\{1,\ldots,r-1\}\), there is some \(t\in T\) such that $$\Delta_{w_r}\cap \left(\bigcup_{i\in T}\Delta_{w_i}\right) = \Delta_{w_r}\cap \Delta_{w_t}.$$ 
    \end{enumerate}
\end{definition}


\iffalse
\tim{Working, may delete

\begin{definition}\label{def:geometrically prunable}
    For any tuple $\bar{i} = [i,j,p,q]$ such that $1<i<j<n$, $i< p\leq j$, $i\leq q<j$, let $Y_{ijpq}$ be the following geometric conditions:
    \[\begin{array}{lcl}
       \bullet\;\; V_1 = \bC^1,\dots, V_{i-1} = \bC^{i-1},  & \hspace{1in} & \bullet\;\;\bC^{i}\subseteq V_{p}, \\
       \bullet\;\; V_{q-1}\subseteq \bC^{j}, \text{ and }  & & \bullet\;\; V_{j} =\bC^{j},\dots, V_{n-1} = \bC^{n-1}
    \end{array}\]
\end{definition}

\begin{definition}    
    Given a locus $Y\subseteq \fl_n$, we say $Y$ is \textbf{geometrically prunable} if for some $0\leq k\leq n$ and all $1\leq r\leq k$ there exists an ordered set of oredered tuples $[i_r,j_r,p_r,q_r] = \bar{i}_r$ such that $Y = \bigcup_{j=1}^k Y_{\bar{i}_j}$ and such that the following two conditions holds for all $2\leq r \leq k$ where $Z_{r} = Y_{\bar{i}_r} \cap \left(\bigcup_{j=1}^{r-1}  Y_{\bar{i}_j}\right)$:
    \begin{enumerate}
    \item[(a)]$\dim(Z_{r}) = \dim(Y_{\bar{i}_r})$ 
    \item[(b)] For every nonempty subset \(T\subseteq\{1,\ldots,r-1\}\), there is some \(t\in T\) such that $$Y_{\bar{i}_r}\cap \left(\bigcup_{j\in T}Y_{\bar{i}_j}\right) = Y_{\bar{i}_r}\cap Y_{\bar{i}_t}.$$ 
    \end{enumerate}
\end{definition}}
\fi

\begin{example}
In the case of $|\calW| = 2$, $\calW = \{v,w\}$ is prunable if and only if the intersection $\Delta_v\cap \Delta_w$ is pure of dimension $\min(\dim\Delta_v,\dim\Delta_w)$. As an example of how \emph{prunable} is not solely a topological condition on the union $\Delta_v\cup \Delta_w$, observe that for $v = 3412$, $w=3241$, $\Delta_v\cap\Delta_w$ is not pure dimensional, so $\{v,w\}$ is not prunable. However, $\Delta_v\cup \Delta_w = \Delta_{3421}$, which is a doubly pruned complex, hence trivially prunable.

\[
\begin{tikzpicture}[scale=0.56, x=1cm, y=1cm,
  vtx/.style={circle, fill=black, inner sep=1.15pt},
  ghost/.style={circle, draw=black!30, fill=white, inner sep=1.0pt},
  ed/.style={line width=0.7pt, black!80}]
  \coordinate (v2) at (0.000,1.300);
  \coordinate (v3) at (-1.220,-0.780);
  \coordinate (v4) at (1.220,-0.780);
  \coordinate (v1) at (0.000,-0.060);
  \fill[faceblue] (v1) -- (v2) -- (v4) -- cycle;
  \fill[faceblue] (v1) -- (v3) -- (v4) -- cycle;
  \draw[ed] (v1) -- (v2);
  \draw[ed] (v1) -- (v3);
  \draw[ed] (v1) -- (v4);
  \draw[ed] (v2) -- (v3);
  \draw[ed] (v2) -- (v4);
  \draw[ed] (v3) -- (v4);
  \node[vtx] at (v1) {};
  \node[above right, inner sep=1.6pt, font=\scriptsize] at (v1) {$1$};
  \node[vtx] at (v2) {};
  \node[above, inner sep=1.6pt, font=\scriptsize] at (v2) {$2$};
  \node[vtx] at (v3) {};
  \node[below left, inner sep=1.6pt, font=\scriptsize] at (v3) {$3$};
  \node[vtx] at (v4) {};
  \node[below right, inner sep=1.6pt, font=\scriptsize] at (v4) {$4$};
  \node at (0,-2) {$\Delta_{3412}$};
\end{tikzpicture}
\qquad
\begin{tikzpicture}[scale=0.56, x=1cm, y=1cm,
  vtx/.style={circle, fill=black, inner sep=1.15pt},
  ghost/.style={circle, draw=black!30, fill=white, inner sep=1.0pt},
  ed/.style={line width=0.7pt, black!80}]
  \coordinate (v2) at (0.000,1.300);
  \coordinate (v3) at (-1.220,-0.780);
  \coordinate (v4) at (1.220,-0.780);
  \coordinate (v1) at (0.000,-0.060);
  \fill[faceblue] (v1) -- (v2) -- (v3) -- cycle;
  \fill[faceblue] (v1) -- (v2) -- (v4) -- cycle;
  \draw[ed] (v1) -- (v2);
  \draw[ed] (v1) -- (v3);
  \draw[ed] (v1) -- (v4);
  \draw[ed] (v2) -- (v3);
  \draw[ed] (v2) -- (v4);
  \node[vtx] at (v1) {};
  \node[above right, inner sep=1.6pt, font=\scriptsize] at (v1) {$1$};
  \node[vtx] at (v2) {};
  \node[above, inner sep=1.6pt, font=\scriptsize] at (v2) {$2$};
  \node[vtx] at (v3) {};
  \node[below left, inner sep=1.6pt, font=\scriptsize] at (v3) {$3$};
  \node[vtx] at (v4) {};
  \node[below right, inner sep=1.6pt, font=\scriptsize] at (v4) {$4$};
    \node at (0,-2) {$\Delta_{3241}$};
\end{tikzpicture}
\quad
\begin{tikzpicture}[scale=0.56, x=1cm, y=1cm,
  vtx/.style={circle, fill=black, inner sep=1.15pt},
  ghost/.style={circle, draw=black!30, fill=white, inner sep=1.0pt},
  ed/.style={line width=0.7pt, black!80}]
  \coordinate (v2) at (0.000,1.300);
  \coordinate (v3) at (-1.220,-0.780);
  \coordinate (v4) at (1.220,-0.780);
  \coordinate (v1) at (0.000,-0.060);
  \fill[faceblue] (v1) -- (v2) -- (v4) -- cycle;
  \draw[ed] (v1) -- (v2);
  \draw[ed] (v1) -- (v3);
  \draw[ed] (v1) -- (v4);
  \draw[ed] (v2) -- (v3);
  \draw[ed] (v2) -- (v4);
  \node[vtx] at (v1) {};
  \node[above right, inner sep=1.6pt, font=\scriptsize] at (v1) {$1$};
  \node[vtx] at (v2) {};
  \node[above, inner sep=1.6pt, font=\scriptsize] at (v2) {$2$};
  \node[vtx] at (v3) {};
  \node[below left, inner sep=1.6pt, font=\scriptsize] at (v3) {$3$};
  \node[vtx] at (v4) {};
  \node[below right, inner sep=1.6pt, font=\scriptsize] at (v4) {$4$};
    \node at (0,-2) {$\Delta_{3412}\cap\Delta_{3241}$};
\end{tikzpicture}
\qquad
\begin{tikzpicture}[scale=0.56, x=1cm, y=1cm,
  vtx/.style={circle, fill=black, inner sep=1.15pt},
  ghost/.style={circle, draw=black!30, fill=white, inner sep=1.0pt},
  ed/.style={line width=0.7pt, black!80}]
  \coordinate (v2) at (0.000,1.300);
  \coordinate (v3) at (-1.220,-0.780);
  \coordinate (v4) at (1.220,-0.780);
  \coordinate (v1) at (0.000,-0.060);
  \fill[faceblue] (v1) -- (v2) -- (v3) -- cycle;
  \fill[faceblue] (v1) -- (v2) -- (v4) -- cycle;
  \fill[faceblue] (v1) -- (v3) -- (v4) -- cycle;
  \draw[ed] (v1) -- (v2);
  \draw[ed] (v1) -- (v3);
  \draw[ed] (v1) -- (v4);
  \draw[ed] (v2) -- (v3);
  \draw[ed] (v2) -- (v4);
  \draw[ed] (v3) -- (v4);
  \node[vtx] at (v1) {};
  \node[above right, inner sep=1.6pt, font=\scriptsize] at (v1) {$1$};
  \node[vtx] at (v2) {};
  \node[above, inner sep=1.6pt, font=\scriptsize] at (v2) {$2$};
  \node[vtx] at (v3) {};
  \node[below left, inner sep=1.6pt, font=\scriptsize] at (v3) {$3$};
  \node[vtx] at (v4) {};
  \node[below right, inner sep=1.6pt, font=\scriptsize] at (v4) {$4$};
    \node at (0,-2) {$\Delta_{3421}$};
\end{tikzpicture}
\]
\end{example}

Along the way, we prove that each simplicial complex $\Delta_w$ for $w\in S_n(\mathcal{P})$ has the following topological properties.



\begin{restatable}[]{thmintro}{MainC}
\label{thm:main-C}
    For each $w\in\sn(\calP),$ the simplicial complex $\Delta_w$ is vertex decomposable, hence shellable and sequentially Cohen-Macaulay. When $\Delta_w$ is pure, it is Cohen-Macaulay. 
\end{restatable}

Moreover, we give simple conditions to determine the purity of a simplicial complex $\Delta_w$ for $w\in S_n(\mathcal P)$.


In \Cref{sec: the combinatorics of doubly pruned simplicial complexes}, we analyze the homological properties of the cohomology ring $R_n/J_w \cong H^*(X_w;\bK)$ and the Stanley--Reisner ring $R_{\Delta_w} \coloneqq S/I_{\Delta_w}$ when $w\in S_n(\calP)$. In particular, we compute the minimal free resolution, the depth and dimension, and the Hilbert series of these rings. The bulk of the work occurs in the case when the squarefree monomial relations ``cross.'' Furthermore, since the elementary symmetric polynomials do not form a regular sequence on $R_{\Delta_w}$, the Hilbert series cannot readily be read off of a Koszul complex computation.

The same failure of freeness is the focus of our last result, which is a more general statement than simply doubly pruned complexes, but does include them as a specialization. After passing to $R_\Delta,$ the relations $I_n$ become canonical. The image of $e_k$ is the sum of $x_F$ over the $k$-element faces of $\Delta,$ the $k$th universal parameter of $\Delta$ denoted $\theta_k$ in the sense of Herzog-Moradi \cite{herzog-moradi} and Adams-Reiner \cite{adams-reiner}, so that $R_{\Delta,n} = R_{\Delta}/(\theta),$ where $\theta = (\theta_1,\ldots,\theta_d).$ Adams and Reiner conjecture \cite[Conj.~6.1]{adams-reiner} that the minimal free resolution of $R_\Delta$ over $\bK[\theta]$ is governed by the barycentric subdivision of $\Delta,$
\[
\dim_\bK
\Tor_r^{\bK[\theta]}(R_\Delta,\bK)_N
=
\sum_{
\substack{
T\subseteq[d] \\
\wt T = N
}
}
\dim_\bK\tilde{H}_{\#T-r-1}(\Sd_T(\Delta);\bK),
\]
where $d = \dim\Delta +1, \wt T = \sum_{k\in T}k,$ and $\Sd_T(\Delta)$ is the rank-selected subcomplex. The conjecture is open in general.

\begin{restatable}[]{thmintro}{MainD}
\label{thm:main-D}
        The Adams-Reiner conjecture (restated in \Cref{eq: Adams-Reiner conj}) holds over any field for every simplicial complex with at most two minimal nonfaces that are of cardinality at least two. In particular, it holds for all doubly pruned complexes, and so for every $\Delta_w$ with $w\in\sn(\calP).$ Moreover, $\Tor_r^{\bK[\theta]}(R_\Delta,\bK) = 0$ for $r\geq 2,$ and $\Tor_1$ vanishes unless the two nonfaces meet. In the latter case,
    \[
    \sum_N
    \dim_\bK \Tor_1^{\bK[\theta]}(R_\Delta,\bK)_N t^N 
    = t^{m-1+\gamma}~[\alpha]_t[\beta]_t[m-2]_t!.
    \]
\end{restatable}

 For $w\in\sn(\calP),$ $\Tor_1$ measures the failure of $H^*(X_w;\bK)$ to be free over $\bK[\theta],$ is nonzero precisely in the crossing case, and has total dimension $(q-i)(j-p)(j-i-1)!.$ Crossing is therefore the failure of Cohen-Macaulayness, the obstruction to the Koszul computation, and the only source of higher $\Tor$ over the parameter ring.

\textbf{Organization.}
This paper is organized as follows. In \Cref{sec: background}, we collect relevant background. In \Cref{sec: ideals in coinvariant rings}, we prove several lemmas on squarefree monomials in $R_n.$ In \Cref{sec: intro to doubly pruned simplicial complexes}, we introduce doubly pruned simplicial complexes and study their properties, including characterizing them by pattern avoidance. In \Cref{sec: proofs of the main theorems}, we prove \Cref{thm:main-A} and \Cref{thm:main-B}, the latter using Koszul cycle arguments. In \Cref{sec: the combinatorics of doubly pruned simplicial complexes}, we prove \Cref{thm:main-C} and \Cref{thm:main-D} and compute the Hilbert series.  In \Cref{sec: example}, we work out the example case where $n=4.$ 

\textbf{Future work.} Squarefree monomials have been studied previously in Schubert geometry in the context of degenerations of matrix Schubert varieties and their coordinate rings, see \cite{bertiger,brion-lakshmibai, hamaker-pechenik-weigandt,klein-diagonal,knutson-lam-speyer-richardson, knutson-miller-grobner}. In contrast, this article characterizes when Schubert cohomology ideals are generated by squarefree monomials. It would be interesting to find parallels between the simplicial complexes in that setting and the doubly pruned complexes introduced in our work.

Product formulas for Poincar\'e polynomials have precedents in Gasharov \cite{gasharov-factoring} and in the fibre bundle structures of Schubert varieties \cite{alland-richmond,richmond-slofstra}, whereas ours come from a free resolution. Closest to Section~6 are the complexes built from lower Bruhat intervals by Bolognini-Sentinelli \cite{bolognini-sentinelli}.

Finally, one question these results raise is whether \Cref{thm:main-B} holds $T$-equivariantly, with squarefree products of the linear forms $x_r-y_s$ in place of squarefree monomials. One implication is immediate upon specializing $y=0,$ and Stanley-Reisner presentations of equivariant cohomology are known in the GKM setting \cite{goertsches-solomadin}, of which $\fl_n$ is an instance.




